problem:  
Let \( O^{2n} \) be a compact, simply connected **torus orbifold** with an effective, isometric \(T^n\)-action such that the orbit space \( Q = O / T^n \) is a manifold with corners combinatorially equivalent to a simple convex polytope.  
Assume that \(O\) is equipped with a Riemannian metric of **nonnegative sectional curvature**.

Weighted projective spaces and, more generally, **toric orbifolds** often inherit nonnegative (even positive) curvature from the standard metric on spheres by quotienting via weighted torus actions. These examples exhibit orbifold singularities but otherwise have rich toric combinatorics.

> **Open problem (Curvature–Singularity Compatibility Problem):**  
> Classify all toric orbifolds \( O^{2n} \) with polytope‑type orbit space that admit a Riemannian metric of nonnegative sectional curvature invariant under the torus action.  
>  
> Specifically, determine which **non‑unimodular characteristic data** \(\lambda : \{\text{facets}(Q)\}\to\mathbb{Z}^n\) can arise from such metrics. Are only certain combinatorial or lattice types compatible with nonnegative curvature?  
>  
> Equivalently: among toric orbifolds \( (Q,\lambda) \), which ones belong to the **metric closure** of the class of nonnegatively curved torus manifolds?

Why is it a "good" problem:  
1. **Corrects a logical flaw and targets genuine unknowns:**  
   Instead of asserting that curvature forbids orbifold singularities — known to be false — the problem seeks to identify which orbifold singularities and lattice‑data types actually *support* nonnegative curvature. No general classification or characterization exists, even in low dimensions.

2. **Conceptual and structural depth:**  
   It asks for the fundamental interface between geometric curvature conditions and discrete combinatorial invariants of toric orbifolds (weights, isotropy lattices). Such a classification would reveal the precise **metric–combinatorial constraints** of curvature within toric topology, bridging continuous and discrete aspects of high symmetry spaces.

3. **Bridges multiple active areas:**  
   The problem connects  
   - **Riemannian and Alexandrov geometry** (nonnegative curvature, orbit space metrics),  
   - **topological transformation group theory** (torus actions and isotropy structures), and  
   - **toric and polyhedral geometry** (characteristic data, unimodularity, weighted projective structures).  
   Its resolution would situate toric topological data within the quantitative framework of curvature geometry.

4. **Simplicity with latent depth:**  
   The core question—“Which toric orbifold singularities are compatible with nonnegative curvature?”—is immediately comprehensible, yet it demands a deep mix of geometric analysis, topology, and polyhedral combinatorics to answer.

5. **Feasible but open:**  
   Known examples (weighted projective spaces, certain Eschenburg or Aloff–Wallach orbifolds) show compatibility in isolated cases, but there is no systematic theory describing what combinatorial data \(\lambda\) allow nonnegative curvature. This makes the problem both well‑defined and currently unsolved.

Thus, this is a “good” problem because it sets a realistic and unanswered goal: the **geometric classification** of nonnegatively curved toric orbifolds and the discovery of the precise **combinatorial constraints** under which curvature and symmetry coexist. It reformulates a previously false rigidity conjecture into a fertile, structurally meaningful open research direction connecting curvature geometry with toric combinatorics.